\part{III gimnazijos klasė}

\section{Sekos ir matematinė indukcija}\label{sec:sekos-ir-matematine-indukcija}

\subsection{Apibrėžtys}

\begin{apibreztis}
Seka yra funkcija $a:\mathbb{N}\to \mathbb{R}$, kurios reikšmė $a(n)$ žymima $a_n$. Skaičius $a_n$ vadinamas $n$-uoju sekos nariu.
\end{apibreztis}

\begin{apibreztis}
Seka $(a_n)$ vadinama aritmetine progresija, jei egzistuoja pastovusis skirtumas $d$, toks kad $a_{n+1}-a_n=d$ visiems $n\in\mathbb{N}$. Ji vadinama geometrine progresija, jei egzistuoja pastovusis vardiklis $q$, toks kad $a_{n+1}=qa_n$ visiems $n\in\mathbb{N}$.
\end{apibreztis}

\begin{apibreztis}
Matematinės indukcijos principas teigia: jei teiginys $P(1)$ teisingas ir iš $P(k)$ teisingumo išplaukia $P(k+1)$ teisingumas kiekvienam $k\in\mathbb{N}$, tai $P(n)$ teisingas visiems $n\in\mathbb{N}$.
\end{apibreztis}

\subsection{Teoremos ir savybės}

\begin{teorema}
Aritmetinės progresijos bendrasis narys ir pirmųjų $n$ narių suma yra
\[
a_n=a_1+(n-1)d,\qquad S_n=\frac{n(a_1+a_n)}{2}.
\]
\emph{Įrodymas.} Iš apibrėžties $a_2=a_1+d$, $a_3=a_1+2d$ ir indukcija gauname $a_n=a_1+(n-1)d$. Sumai parašome dvi eilutes:
\[
S_n=a_1+a_2+\cdots+a_n,\qquad S_n=a_n+a_{n-1}+\cdots+a_1.
\]
Sudėję gauname $2S_n=(a_1+a_n)+\cdots+(a_n+a_1)=n(a_1+a_n)$.
\end{teorema}

\begin{teorema}
Geometrinės progresijos bendrasis narys ir suma, kai $q\ne 1$, yra
\[
a_n=a_1q^{n-1},\qquad S_n=a_1\frac{q^n-1}{q-1}.
\]
\emph{Įrodymas.} Bendrojo nario formulė gaunama kartotinai dauginant iš $q$. Sumai taikome
\[
qS_n=a_1q+a_1q^2+\cdots+a_1q^n.
\]
Atėmę $S_n$ iš $qS_n$, gauname $(q-1)S_n=a_1(q^n-1)$.
\end{teorema}

\begin{savybe}
Jei seka didėjanti ir aprėžta iš viršaus, ji turi baigtinę ribą. Jei seka mažėjanti ir aprėžta iš apačios, ji taip pat turi baigtinę ribą.
\end{savybe}

\subsection{Taisyklės ir algoritmai}

Norint taikyti indukciją, pirmiausia patikrinamas pradinis atvejis. Tada aiškiai suformuluojama indukcijos prielaida $P(k)$ ir, naudojant tik ją bei jau žinomas tapatybes, įrodomas $P(k+1)$.

Sekos rekursinei formulei tirti dažnai ieškoma pastoviojo taško $L$, sprendžiant lygtį $L=f(L)$, tada įrodoma monotonija ir aprėžtumas. Tik po to leidžiama pereiti prie ribos.

\subsection{Iliustruojantys pavyzdžiai}

\textbf{Pavyzdys 1.} Įrodykime formulę
\[
1+2+\cdots+n=\frac{n(n+1)}{2}.
\]
Kai $n=1$, abi pusės lygios $1$. Tarkime, kad formulei teisinga, kai $n=k$, t. y.
\[
1+\cdots+k=\frac{k(k+1)}{2}.
\]
Tada
\[
1+\cdots+k+(k+1)=\frac{k(k+1)}{2}+k+1=\frac{(k+1)(k+2)}{2}.
\]
Tai yra formulė, kai $n=k+1$. Pagal indukciją ji teisinga visiems $n$. \textbf{Atsakymas: formulė įrodyta visiems $n\in\mathbb{N}$.}

\textbf{Pavyzdys 2.} Raskime aritmetinės progresijos sumą, jei $a_1=7$, $d=3$, $n=20$. Turime
\[
a_{20}=7+19\cdot 3=64,
\]
todėl
\[
S_{20}=\frac{20(7+64)}{2}=10\cdot 71=710.
\]
\textbf{Atsakymas: $710$.}

\textbf{Pavyzdys 3.} Rekursiškai apibrėžta seka $a_1=1$, $a_{n+1}=\sqrt{2+a_n}$. Raskime jos ribą. Pirmiausia įrodome, kad $a_n<2$: jei $a_k<2$, tai $a_{k+1}<\sqrt4=2$, o $a_1<2$. Be to,
\[
a_{n+1}\ge a_n \iff \sqrt{2+a_n}\ge a_n \iff 2+a_n\ge a_n^2 \iff (a_n-2)(a_n+1)\le0,
\]
kas teisinga, nes $-1<a_n<2$. Seka didėjanti ir aprėžta, todėl turi ribą $L$. Tada
\[
L=\sqrt{2+L},\qquad L^2-L-2=0,\qquad L\in\{2,-1\}.
\]
Kadangi $a_n>0$, gauname $L=2$. \textbf{Atsakymas: $\lim a_n=2$.}

\subsection{Ryšys su kitomis temomis}

Sekos remiasi funkcijų samprata iš \ref{sec:funkcijos-ir-ju-savybes}, o ribos natūraliai pereina į \ref{sec:ribos-ir-tolydumas}. Indukcija yra pagrindinis įrodymo metodas kombinatorikoje iš \ref{sec:tikimybe-ir-kombinatorika} ir algebrinėse formulėse, naudojamose \ref{sec:laipsnines-rodiklines-ir-logaritmines-funkcijos}.

\section{Laipsninės, rodiklinės ir logaritminės funkcijos}\label{sec:laipsnines-rodiklines-ir-logaritmines-funkcijos}

\subsection{Apibrėžtys}

\begin{apibreztis}
Laipsninė funkcija yra funkcija $f(x)=x^\alpha$, kur $\alpha\in\mathbb{R}$, o apibrėžimo sritis parenkama taip, kad reiškinys turėtų realią prasmę.
\end{apibreztis}

\begin{apibreztis}
Rodiklinė funkcija yra $f(x)=a^x$, kur $a>0$, $a\ne1$. Logaritminė funkcija yra $f(x)=\log_a x$, kur $a>0$, $a\ne1$, $x>0$, ir ji apibrėžiama sąlyga
\[
\log_a x=y \iff a^y=x.
\]
\end{apibreztis}

\begin{pastaba}
Funkcijos $a^x$ ir $\log_a x$ yra viena kitai atvirkštinės: $a^{\log_a x}=x$ ir $\log_a(a^x)=x$.
\end{pastaba}

\subsection{Teoremos ir savybės}

\begin{teorema}
Jei $a>0$, $a\ne1$, tai $a^x$ yra griežtai didėjanti, kai $a>1$, ir griežtai mažėjanti, kai $0<a<1$.
\emph{Įrodymas.} Kai $a>1$ ir $x<y$, turime $y-x>0$, todėl $a^{y}=a^x a^{y-x}>a^x$. Kai $0<a<1$, rašome $a=1/b$, kur $b>1$, ir gauname $a^x=b^{-x}$; didėjant $x$, mažėja $-x$, todėl $b^{-x}$ mažėja.
\end{teorema}

\begin{savybe}
Logaritmams galioja
\[
\log_a(xy)=\log_a x+\log_a y,\quad
\log_a\frac{x}{y}=\log_a x-\log_a y,\quad
\log_a(x^r)=r\log_a x.
\]
\end{savybe}

\begin{teorema}
Pagrindo keitimo formulė:
\[
\log_a x=\frac{\log_b x}{\log_b a}.
\]
\emph{Įrodymas.} Tegul $y=\log_a x$. Tada $a^y=x$. Pritaikę $\log_b$, gauname $y\log_b a=\log_b x$, todėl $y=\dfrac{\log_b x}{\log_b a}$.
\end{teorema}

\subsection{Taisyklės ir algoritmai}

Sprendžiant rodiklines lygtis, siekiama suvesti abi puses į tą patį pagrindą arba taikyti logaritmavimą. Sprendžiant logaritmines lygtis, pirmiausia nustatoma apibrėžimo sritis, tada naudojamos logaritmų savybės, o gauti sprendiniai patikrinami pradinėje srityje.

Nelygybėse būtina atsižvelgti į monotoniškumą: jei $a>1$, tai $\log_a x<\log_a y \iff x<y$, o jei $0<a<1$, nelygybės ženklas apsiverčia.

\subsection{Iliustruojantys pavyzdžiai}

\textbf{Pavyzdys 1.} Išspręskime lygtį $3^{2x-1}=27^{x-2}$. Kadangi $27=3^3$, turime
\[
3^{2x-1}=3^{3x-6}.
\]
Funkcija $3^t$ griežtai didėjanti, todėl
\[
2x-1=3x-6,\qquad x=5.
\]
\textbf{Atsakymas: $x=5$.}

\textbf{Pavyzdys 2.} Išspręskime $\log_2(x-1)+\log_2(x+3)=3$. Apibrėžimo sritis: $x>1$. Tada
\[
\log_2((x-1)(x+3))=3,
\]
\[
(x-1)(x+3)=8,\qquad x^2+2x-3=8,\qquad x^2+2x-11=0.
\]
Gauname
\[
x=-1\pm 2\sqrt3.
\]
Sąlygą $x>1$ tenkina tik $x=-1+2\sqrt3$. \textbf{Atsakymas: $x=-1+2\sqrt3$.}

\textbf{Pavyzdys 3.} Raskime funkcijos $f(x)=\log_3(9-x^2)$ apibrėžimo sritį ir reikšmių sritį. Reikia
\[
9-x^2>0 \iff -3<x<3.
\]
Kadangi $9-x^2\in(0,9]$, o $\log_3$ didėjanti, tai
\[
f(x)\in(-\infty,\log_3 9]=(-\infty,2].
\]
\textbf{Atsakymas: $D(f)=(-3,3)$, $R(f)=(-\infty,2]$.}

\subsection{Ryšys su kitomis temomis}

Ši tema pratęsia \ref{sec:laipsniai-saknys-ir-logaritmu-pradzia} ir remiasi realiųjų skaičių intervalais iš \ref{sec:realieji-skaiciai-ir-intervalai}. Rodiklinės funkcijos modeliuoja augimą sekose iš \ref{sec:sekos-ir-matematine-indukcija}, o logaritminis diferencijavimas naudojamas \ref{sec:isvestine-ir-diferencijavimas}.

\section{Trigonometrinės funkcijos}\label{sec:trigonometrines-funkcijos}

\subsection{Apibrėžtys}

\begin{apibreztis}
Vienetiniame apskritime kampui $t$ atitinkantis taškas turi koordinates $(\cos t,\sin t)$. Taip apibrėžiamos funkcijos $\sin t$ ir $\cos t$ visiems realiesiems $t$.
\end{apibreztis}

\begin{apibreztis}
Funkcijos $\tan t$ ir $\cot t$ apibrėžiamos
\[
\tan t=\frac{\sin t}{\cos t}\quad(\cos t\ne0),\qquad
\cot t=\frac{\cos t}{\sin t}\quad(\sin t\ne0).
\]
\end{apibreztis}

\begin{apibreztis}
Funkcija $f$ vadinama periodine, jei egzistuoja $T>0$, toks kad $f(x+T)=f(x)$ visiems apibrėžimo srities $x$. Mažiausias toks $T$, jei jis egzistuoja, vadinamas pagrindiniu periodu.
\end{apibreztis}

\subsection{Teoremos ir savybės}

\begin{savybe}
Galioja pagrindinė tapatybė
\[
\sin^2 x+\cos^2 x=1.
\]
Ji išplaukia iš vienetinio apskritimo lygties $X^2+Y^2=1$.
\end{savybe}

\begin{savybe}
Funkcijos $\sin x$ ir $\cos x$ turi periodą $2\pi$, o $\tan x$ ir $\cot x$ turi periodą $\pi$.
\end{savybe}

\begin{teorema}
Sudėties formulės:
\[
\sin(x+y)=\sin x\cos y+\cos x\sin y,
\]
\[
\cos(x+y)=\cos x\cos y-\sin x\sin y.
\]
\emph{Įrodymas.} Sukimas kampu $t$ plokštumoje aprašomas matrica
\[
R_t=\begin{pmatrix}\cos t&-\sin t\\ \sin t&\cos t\end{pmatrix}.
\]
Kadangi $R_xR_y=R_{x+y}$, sulyginus matricų elementus gaunamos abi formulės.
\end{teorema}

\subsection{Taisyklės ir algoritmai}

Grafikams braižyti pirmiausia nustatomas periodas, amplitudė ir poslinkiai. Funkcijai $A\sin(bx+c)+d$ amplitudė yra $|A|$, periodas $\frac{2\pi}{|b|}$, vidurio linija $y=d$.

Skaičiuojant tikslias reikšmes, kampai redukuojami į pirmąjį ketvirtį ir naudojamos reikšmės kampams $0$, $\frac{\pi}{6}$, $\frac{\pi}{4}$, $\frac{\pi}{3}$, $\frac{\pi}{2}$.

\subsection{Iliustruojantys pavyzdžiai}

\textbf{Pavyzdys 1.} Raskime $f(x)=3\sin(2x-\pi)+1$ amplitudę, periodą ir vidurio liniją. Čia $A=3$, $b=2$, $d=1$, todėl amplitudė $3$, periodas
\[
T=\frac{2\pi}{2}=\pi,
\]
o vidurio linija $y=1$. \textbf{Atsakymas: amplitudė $3$, periodas $\pi$, vidurio linija $y=1$.}

\textbf{Pavyzdys 2.} Apskaičiuokime $\sin\frac{7\pi}{6}$ ir $\cos\frac{7\pi}{6}$. Kampas $\frac{7\pi}{6}=\pi+\frac{\pi}{6}$ yra trečiajame ketvirtyje, todėl abi reikšmės neigiamos:
\[
\sin\frac{7\pi}{6}=-\sin\frac{\pi}{6}=-\frac12,\qquad
\cos\frac{7\pi}{6}=-\cos\frac{\pi}{6}=-\frac{\sqrt3}{2}.
\]
\textbf{Atsakymas: $\sin\frac{7\pi}{6}=-\frac12$, $\cos\frac{7\pi}{6}=-\frac{\sqrt3}{2}$.}

\textbf{Pavyzdys 3.} Įrodykime, kad $\sin 2x=2\sin x\cos x$. Iš sudėties formulės
\[
\sin(x+x)=\sin x\cos x+\cos x\sin x=2\sin x\cos x.
\]
Kadangi $x+x=2x$, gauname norimą tapatybę. \textbf{Atsakymas: $\sin 2x=2\sin x\cos x$.}

\subsection{Ryšys su kitomis temomis}

Trigonometrinės funkcijos pratęsia trigonometriją stačiajame trikampyje iš \ref{sec:trigonometrija-staciajame-trikampyje} ir apskritimo geometriją iš \ref{sec:apskritimo-geometrija}. Jų grafikai ir transformacijos remiasi \ref{sec:funkcijos-ir-ju-savybes}, o išvestinės nagrinėjamos \ref{sec:isvestine-ir-diferencijavimas}.

\section{Trigonometrinės lygtys ir tapatybės}\label{sec:trigonometrines-lygtys-ir-tapatybes}

\subsection{Apibrėžtys}

\begin{apibreztis}
Trigonometrinė lygtis yra lygtis, kurioje nežinomasis yra trigonometrinės funkcijos argumente arba reiškinyje su trigonometrinėmis funkcijomis.
\end{apibreztis}

\begin{apibreztis}
Trigonometrinė tapatybė yra lygybė, teisinga visoms argumento reikšmėms, kurioms abi pusės apibrėžtos.
\end{apibreztis}

\begin{pastaba}
Bendrieji sprendiniai rašomi su $k\in\mathbb{Z}$, nes trigonometrinės funkcijos yra periodinės.
\end{pastaba}

\subsection{Teoremos ir savybės}

\begin{savybe}
Pagrindinių lygčių sprendiniai:
\[
\sin x=a \iff x=(-1)^k\arcsin a+k\pi,\quad |a|\le1,
\]
\[
\cos x=a \iff x=\pm\arccos a+2k\pi,\quad |a|\le1,
\]
\[
\tan x=a \iff x=\arctan a+k\pi.
\]
\end{savybe}

\begin{savybe}
Iš pagrindinės tapatybės gaunama
\[
1+\tan^2 x=\frac{1}{\cos^2 x},\qquad 1+\cot^2 x=\frac{1}{\sin^2 x}.
\]
\end{savybe}

\begin{teorema}
Sandaugos nulio principas trigonometrinėms lygtims taikomas taip pat kaip algebrinėms: jei $A(x)B(x)=0$, tai $A(x)=0$ arba $B(x)=0$.
\emph{Įrodymas.} Kiekvienai leistinai $x$ reikšmei $A(x)$ ir $B(x)$ yra realieji skaičiai. Realiųjų skaičių sandauga lygi nuliui tada ir tik tada, kai bent vienas dauginamasis lygus nuliui.
\end{teorema}

\subsection{Taisyklės ir algoritmai}

Sprendžiant lygtį pirmiausia nustatoma apibrėžimo sritis. Tada lygtis pertvarkoma iki pagrindinių lygčių arba sandaugos. Jei buvo dalyta iš reiškinio su nežinomuoju, būtina patikrinti, ar neprarasti sprendiniai.

Tapatybei įrodyti dažniausiai pertvarkoma sudėtingesnė pusė, naudojamos $\sin^2x+\cos^2x=1$, sudėties, dvigubo kampo ir redukcijos formulės.

\subsection{Iliustruojantys pavyzdžiai}

\textbf{Pavyzdys 1.} Išspręskime $2\sin x-1=0$. Gauname $\sin x=\frac12$, todėl
\[
x=\frac{\pi}{6}+2k\pi \quad \text{arba}\quad x=\frac{5\pi}{6}+2k\pi,\qquad k\in\mathbb{Z}.
\]
\textbf{Atsakymas: $x=\frac{\pi}{6}+2k\pi$ arba $x=\frac{5\pi}{6}+2k\pi$, $k\in\mathbb{Z}$.}

\textbf{Pavyzdys 2.} Išspręskime $2\cos^2x-3\cos x+1=0$. Pažymėję $u=\cos x$, gauname
\[
2u^2-3u+1=0,\qquad (2u-1)(u-1)=0.
\]
Taigi $\cos x=\frac12$ arba $\cos x=1$. Todėl
\[
x=\pm\frac{\pi}{3}+2k\pi \quad \text{arba}\quad x=2k\pi.
\]
\textbf{Atsakymas: $x=\pm\frac{\pi}{3}+2k\pi$ arba $x=2k\pi$, $k\in\mathbb{Z}$.}

\textbf{Pavyzdys 3.} Įrodykime tapatybę $\dfrac{1-\cos 2x}{\sin 2x}=\tan x$, kai abi pusės apibrėžtos. Naudojame formules
\[
1-\cos 2x=1-(\cos^2x-\sin^2x)=2\sin^2x,
\]
\[
\sin 2x=2\sin x\cos x.
\]
Todėl
\[
\frac{1-\cos 2x}{\sin 2x}=\frac{2\sin^2x}{2\sin x\cos x}=\frac{\sin x}{\cos x}=\tan x.
\]
\textbf{Atsakymas: tapatybė įrodyta jos apibrėžimo srityje.}

\subsection{Ryšys su kitomis temomis}

Ši tema remiasi \ref{sec:trigonometrines-funkcijos} ir algebrinių lygčių technika iš \ref{sec:kvadratines-lygtys-ir-nelygybes}. Tapatybės reikalingos diferencijuojant trigonometrines funkcijas \ref{sec:isvestine-ir-diferencijavimas} ir sprendžiant geometrinius uždavinius iš \ref{sec:analitine-geometrija-ir-vektoriai}.

\section{Ribos ir tolydumas}\label{sec:ribos-ir-tolydumas}

\subsection{Apibrėžtys}

\begin{apibreztis}
Skaičius $L$ vadinamas funkcijos $f$ riba taške $a$, jei kiekvienai $\varepsilon>0$ egzistuoja $\delta>0$, kad iš $0<|x-a|<\delta$ išplaukia $|f(x)-L|<\varepsilon$. Rašoma $\lim_{x\to a}f(x)=L$.
\end{apibreztis}

\begin{apibreztis}
Funkcija $f$ vadinama tolydžia taške $a$, jei $f(a)$ apibrėžta ir
\[
\lim_{x\to a}f(x)=f(a).
\]
\end{apibreztis}

\begin{apibreztis}
Vienpusės ribos $\lim_{x\to a-}f(x)$ ir $\lim_{x\to a+}f(x)$ apibūdina funkcijos elgesį artėjant prie $a$ atitinkamai iš kairės ir iš dešinės.
\end{apibreztis}

\subsection{Teoremos ir savybės}

\begin{teorema}
Jei egzistuoja $\lim f(x)=A$ ir $\lim g(x)=B$, tai egzistuoja sumos, skirtumo, sandaugos ir, kai $B\ne0$, dalmens ribos:
\[
\lim(f+g)=A+B,\quad \lim(fg)=AB,\quad \lim\frac{f}{g}=\frac{A}{B}.
\]
\emph{Įrodymas.} Sumai taikome trikampio nelygybę:
\[
|(f+g)-(A+B)|\le |f-A|+|g-B|.
\]
Parinkus kiekvienai daliai po $\varepsilon/2$, suma mažesnė už $\varepsilon$. Sandaugai rašome $fg-AB=f(g-B)+B(f-A)$ ir naudojame faktą, kad ribą turinti funkcija šalia taško yra aprėžta. Dalmuo gaunamas iš sandaugos ir $\lim(1/g)=1/B$.
\end{teorema}

\begin{teorema}
Jei $f$ ir $g$ tolydžios taške $a$, tai $f+g$, $fg$ ir, kai $g(a)\ne0$, $\frac{f}{g}$ yra tolydžios taške $a$.
\emph{Įrodymas.} Tai tiesiogiai išplaukia iš ribų algebros ir apibrėžties $\lim f(x)=f(a)$, $\lim g(x)=g(a)$.
\end{teorema}

\begin{savybe}
Polinomai tolydūs visoje $\mathbb{R}$, racionaliosios funkcijos tolydžios ten, kur vardiklis nelygus nuliui, o $\sin x$, $\cos x$, $a^x$ ir $\log_a x$ tolydžios savo apibrėžimo srityse.
\end{savybe}

\subsection{Taisyklės ir algoritmai}

Ribai skaičiuoti pirmiausia bandoma tiesiogiai įstatyti argumentą. Jei gaunama neapibrėžtis, reiškinys pertvarkomas: išskaidomas dauginamaisiais, racionalizuojamas, suvedamas į žinomas ribas.

Tolydumui tikrinti taške reikia trijų sąlygų: $f(a)$ egzistuoja, $\lim_{x\to a}f(x)$ egzistuoja, ir abi reikšmės sutampa.

\subsection{Iliustruojantys pavyzdžiai}

\textbf{Pavyzdys 1.} Apskaičiuokime
\[
\lim_{x\to2}\frac{x^2-4}{x-2}.
\]
Kai $x\ne2$,
\[
\frac{x^2-4}{x-2}=\frac{(x-2)(x+2)}{x-2}=x+2.
\]
Todėl
\[
\lim_{x\to2}\frac{x^2-4}{x-2}=\lim_{x\to2}(x+2)=4.
\]
\textbf{Atsakymas: $4$.}

\textbf{Pavyzdys 2.} Nustatykime, ar funkcija
\[
f(x)=\begin{cases}2x+1,&x<1,\\ x^2+2,&x\ge1\end{cases}
\]
tolydi taške $x=1$. Kairioji riba yra $2\cdot1+1=3$, dešinioji riba ir funkcijos reikšmė yra $1^2+2=3$. Vadinasi,
\[
\lim_{x\to1}f(x)=3=f(1).
\]
\textbf{Atsakymas: funkcija tolydi taške $x=1$.}

\textbf{Pavyzdys 3.} Apskaičiuokime
\[
\lim_{x\to0}\frac{\sqrt{1+x}-1}{x}.
\]
Racionalizuojame:
\[
\frac{\sqrt{1+x}-1}{x}\cdot\frac{\sqrt{1+x}+1}{\sqrt{1+x}+1}
=\frac{1+x-1}{x(\sqrt{1+x}+1)}=\frac{1}{\sqrt{1+x}+1}.
\]
Todėl riba lygi $\frac12$. \textbf{Atsakymas: $\frac12$.}

\subsection{Ryšys su kitomis temomis}

Ribos pratęsia sekų analizę iš \ref{sec:sekos-ir-matematine-indukcija} ir yra būtinas pagrindas išvestinei \ref{sec:isvestine-ir-diferencijavimas}. Tolydumas leidžia griežtai kalbėti apie funkcijų grafikus iš \ref{sec:funkcijos-ir-ju-savybes}.

\section{Išvestinė ir diferencijavimas}\label{sec:isvestine-ir-diferencijavimas}

\subsection{Apibrėžtys}

\begin{apibreztis}
Funkcijos $f$ išvestinė taške $a$ yra riba
\[
f'(a)=\lim_{h\to0}\frac{f(a+h)-f(a)}{h},
\]
jei ši riba egzistuoja.
\end{apibreztis}

\begin{apibreztis}
Funkcija vadinama diferencijuojama intervale, jei ji turi išvestinę kiekviename to intervalo taške.
\end{apibreztis}

\begin{pastaba}
Išvestinė $f'(a)$ yra liestinės krypties koeficientas funkcijos grafiko taške $(a,f(a))$ ir momentinis kitimo greitis.
\end{pastaba}

\subsection{Teoremos ir savybės}

\begin{teorema}
Jei $f$ diferencijuojama taške $a$, tai $f$ tolydi taške $a$.
\emph{Įrodymas.} Turime
\[
f(x)-f(a)=\frac{f(x)-f(a)}{x-a}(x-a).
\]
Kai $x\to a$, pirmasis daugiklis artėja prie $f'(a)$, o antrasis prie $0$, todėl $f(x)-f(a)\to0$.
\end{teorema}

\begin{savybe}
Pagrindinės taisyklės:
\[
(cf)'=cf',\qquad (f+g)'=f'+g',\qquad (fg)'=f'g+fg',
\]
\[
\left(\frac{f}{g}\right)'=\frac{f'g-fg'}{g^2},\quad g\ne0.
\]
\end{savybe}

\begin{teorema}
Grandinės taisyklė: jei $g$ diferencijuojama taške $x$, o $f$ diferencijuojama taške $g(x)$, tai
\[
(f\circ g)'(x)=f'(g(x))g'(x).
\]
\emph{Įrodymas.} Pokytį $f(g(x+h))-f(g(x))$ nagrinėjame per tarpinį pokytį $\Delta g=g(x+h)-g(x)$. Diferencijuojamumas reiškia $f(g+\Delta g)-f(g)=f'(g)\Delta g+r(\Delta g)$, kur $r(\Delta g)/\Delta g\to0$. Padalijus iš $h$ ir leidus $h\to0$, gaunama formulė.
\end{teorema}

\subsection{Taisyklės ir algoritmai}

Diferencijuojant sudėtinę funkciją, iš pradžių atpažįstama išorinė ir vidinė funkcijos dalis. Diferencijuojama išorinė funkcija, jos argumentą paliekant vidinį, ir dauginama iš vidinės funkcijos išvestinės.

Naudingos formulės:
\[
(x^n)'=nx^{n-1},\quad (e^x)'=e^x,\quad (\ln x)'=\frac1x,\quad (\sin x)'=\cos x,\quad (\cos x)'=-\sin x.
\]

\subsection{Iliustruojantys pavyzdžiai}

\textbf{Pavyzdys 1.} Iš apibrėžties raskime $f(x)=x^2$ išvestinę taške $a$. Turime
\[
f'(a)=\lim_{h\to0}\frac{(a+h)^2-a^2}{h}
=\lim_{h\to0}\frac{2ah+h^2}{h}
=\lim_{h\to0}(2a+h)=2a.
\]
\textbf{Atsakymas: $f'(a)=2a$.}

\textbf{Pavyzdys 2.} Diferencijuokime $y=(3x^2-1)^5$. Čia išorinė funkcija $u^5$, vidinė $u=3x^2-1$. Todėl
\[
y'=5(3x^2-1)^4\cdot6x=30x(3x^2-1)^4.
\]
\textbf{Atsakymas: $y'=30x(3x^2-1)^4$.}

\textbf{Pavyzdys 3.} Raskime liestinės lygtį funkcijai $f(x)=\ln x$ taške $x=1$. Turime $f(1)=0$, $f'(x)=1/x$, todėl $f'(1)=1$. Liestinė:
\[
y-0=1(x-1).
\]
\textbf{Atsakymas: $y=x-1$.}

\subsection{Ryšys su kitomis temomis}

Išvestinė remiasi ribomis iš \ref{sec:ribos-ir-tolydumas}. Ji taikoma funkcijų tyrimui \ref{sec:isvestines-taikymai}, logaritminėms ir rodiklinėms funkcijoms iš \ref{sec:laipsnines-rodiklines-ir-logaritmines-funkcijos} bei trigonometrinėms funkcijoms iš \ref{sec:trigonometrines-funkcijos}.

\section{Išvestinės taikymai}\label{sec:isvestines-taikymai}

\subsection{Apibrėžtys}

\begin{apibreztis}
Taškas $c$ vadinamas kritiniu funkcijos $f$ tašku, jei $f'(c)=0$ arba $f'(c)$ neegzistuoja, bet $f(c)$ apibrėžta.
\end{apibreztis}

\begin{apibreztis}
Funkcija turi lokalųjį maksimumą taške $c$, jei tam tikroje $c$ aplinkoje $f(x)\le f(c)$. Ji turi lokalųjį minimumą, jei $f(x)\ge f(c)$.
\end{apibreztis}

\begin{apibreztis}
Funkcijos grafiko įgaubtumas apibūdinamas antrosios išvestinės ženklu: kai $f''>0$, grafikas išgaubtas aukštyn, kai $f''<0$, išgaubtas žemyn.
\end{apibreztis}

\subsection{Teoremos ir savybės}

\begin{teorema}
Jei $f$ diferencijuojama intervale ir $f'(x)>0$ visame intervale, tai $f$ tame intervale didėja; jei $f'(x)<0$, tai mažėja.
\emph{Įrodymas.} Jei $x_1<x_2$, pagal Lagranžo vidutinės reikšmės teoremą egzistuoja $c\in(x_1,x_2)$, kad
\[
f(x_2)-f(x_1)=f'(c)(x_2-x_1).
\]
Kai $f'(c)>0$, skirtumas teigiamas; kai $f'(c)<0$, neigiamas.
\end{teorema}

\begin{savybe}
Jei $f'$ keičia ženklą iš $+$ į $-$ taške $c$, funkcija turi lokalųjį maksimumą. Jei $f'$ keičia ženklą iš $-$ į $+$, funkcija turi lokalųjį minimumą.
\end{savybe}

\begin{savybe}
Jei $f''(c)>0$ ir $f'(c)=0$, tai $c$ yra lokaliojo minimumo taškas. Jei $f''(c)<0$ ir $f'(c)=0$, tai $c$ yra lokaliojo maksimumo taškas.
\end{savybe}

\subsection{Taisyklės ir algoritmai}

Funkcijai tirti randama apibrėžimo sritis, išvestinė, kritiniai taškai, išvestinės ženklų intervalai, ekstremumai, prireikus antroji išvestinė ir ribinis elgesys. Optimizavimo uždavinyje pirmiausia sudaroma vieno kintamojo funkcija ir nustatomas leistinas intervalas.

\subsection{Iliustruojantys pavyzdžiai}

\textbf{Pavyzdys 1.} Ištirkime $f(x)=x^3-3x^2+2$ didėjimą ir ekstremumus. Išvestinė
\[
f'(x)=3x^2-6x=3x(x-2).
\]
Ženklai: $f'>0$, kai $x<0$ arba $x>2$; $f'<0$, kai $0<x<2$. Todėl ties $x=0$ yra lokalus maksimumas, ties $x=2$ lokalus minimumas. Reikšmės:
\[
f(0)=2,\qquad f(2)=8-12+2=-2.
\]
\textbf{Atsakymas: didėja $(-\infty,0)$ ir $(2,\infty)$, mažėja $(0,2)$, maksimumas $2$, minimumas $-2$.}

\textbf{Pavyzdys 2.} Raskime stačiakampio didžiausią plotą, jei jo perimetras $40$. Tegul kraštinės $x$ ir $y$. Tada $2x+2y=40$, todėl $y=20-x$. Plotas
\[
A(x)=x(20-x)=20x-x^2,\qquad 0<x<20.
\]
\[
A'(x)=20-2x.
\]
Kritinis taškas $x=10$, o $A''(x)=-2<0$, todėl tai maksimumas. Tada $y=10$, $A=100$. \textbf{Atsakymas: didžiausias plotas $100$, kai stačiakampis yra $10\times10$.}

\textbf{Pavyzdys 3.} Raskime funkcijos $f(x)=x+\frac4x$, $x>0$, mažiausią reikšmę. Išvestinė
\[
f'(x)=1-\frac4{x^2}.
\]
Lygybė $f'(x)=0$ duoda $x=2$. Kai $0<x<2$, $f'<0$, kai $x>2$, $f'>0$, todėl $x=2$ yra minimumas. 
\[
f(2)=2+\frac42=4.
\]
\textbf{Atsakymas: mažiausia reikšmė $4$.}

\subsection{Ryšys su kitomis temomis}

Išvestinės taikymai tęsia \ref{sec:isvestine-ir-diferencijavimas} ir leidžia sistemiškai tirti funkcijas iš \ref{sec:laipsnines-rodiklines-ir-logaritmines-funkcijos} bei \ref{sec:trigonometrines-funkcijos}. Optimizavimas siejasi su modeliavimu iš \ref{sec:optimizavimas-ir-modeliavimas}.

\section{Analitinė geometrija ir vektoriai}\label{sec:analitine-geometrija-ir-vektoriai}

\subsection{Apibrėžtys}

\begin{apibreztis}
Vektorius plokštumoje yra kryptį ir ilgį turintis dydis. Jei $A(x_1,y_1)$ ir $B(x_2,y_2)$, tai
\[
\overrightarrow{AB}=(x_2-x_1,\ y_2-y_1).
\]
\end{apibreztis}

\begin{apibreztis}
Vektorių $\vec a=(a_1,a_2)$ ir $\vec b=(b_1,b_2)$ skaliarinė sandauga yra
\[
\vec a\cdot \vec b=a_1b_1+a_2b_2=|\vec a||\vec b|\cos\theta,
\]
kur $\theta$ yra kampas tarp vektorių.
\end{apibreztis}

\begin{apibreztis}
Tiesės bendroji lygtis plokštumoje yra $Ax+By+C=0$, kur $(A,B)\ne(0,0)$. Vektorius $(A,B)$ yra tiesės normalės vektorius.
\end{apibreztis}

\subsection{Teoremos ir savybės}

\begin{savybe}
Atstumas tarp taškų $A(x_1,y_1)$ ir $B(x_2,y_2)$ yra
\[
AB=\sqrt{(x_2-x_1)^2+(y_2-y_1)^2}.
\]
\end{savybe}

\begin{savybe}
Vektoriai statmeni tada ir tik tada, kai jų skaliarinė sandauga lygi nuliui.
\end{savybe}

\begin{teorema}
Atstumas nuo taško $P(x_0,y_0)$ iki tiesės $Ax+By+C=0$ yra
\[
d=\frac{|Ax_0+By_0+C|}{\sqrt{A^2+B^2}}.
\]
\emph{Įrodymas.} Tiesės normalė yra $\vec n=(A,B)$. Pasirinkime tašką $Q$ tiesėje. Atstumas nuo $P$ iki tiesės yra vektoriaus $\overrightarrow{QP}$ projekcijos į normalės kryptį modulis:
\[
d=\left|\frac{\overrightarrow{QP}\cdot \vec n}{|\vec n|}\right|.
\]
Kadangi $Q$ tenkina $Ax_Q+By_Q+C=0$, skaitiklis lygus $Ax_0+By_0+C$.
\end{teorema}

\subsection{Taisyklės ir algoritmai}

Tiesei per du taškus rasti galima sudaryti krypties vektorių, tada parinkti statmeną normalę. Kampui tarp tiesių rasti naudojamas kampas tarp jų krypties arba normalės vektorių.

Apskritimo lygtis su centru $(a,b)$ ir spinduliu $r$ yra
\[
(x-a)^2+(y-b)^2=r^2.
\]

\subsection{Iliustruojantys pavyzdžiai}

\textbf{Pavyzdys 1.} Raskime tiesės per $A(1,2)$ ir $B(5,4)$ lygtį. Krypties vektorius
\[
\overrightarrow{AB}=(4,2),
\]
tad normalė gali būti $(2,-4)$ arba supaprastinta $(1,-2)$. Lygtis:
\[
1(x-1)-2(y-2)=0,
\]
\[
x-1-2y+4=0,\qquad x-2y+3=0.
\]
\textbf{Atsakymas: $x-2y+3=0$.}

\textbf{Pavyzdys 2.} Apskaičiuokime kampą tarp $\vec a=(1,\sqrt3)$ ir $\vec b=(2,0)$. Turime
\[
\vec a\cdot\vec b=2,\quad |\vec a|=2,\quad |\vec b|=2.
\]
Todėl
\[
\cos\theta=\frac{2}{2\cdot2}=\frac12,
\]
vadinasi $\theta=\frac{\pi}{3}$. \textbf{Atsakymas: $\frac{\pi}{3}$.}

\textbf{Pavyzdys 3.} Raskime atstumą nuo taško $P(3,-1)$ iki tiesės $2x-y+4=0$. Taikome formulę:
\[
d=\frac{|2\cdot3-(-1)+4|}{\sqrt{2^2+(-1)^2}}
=\frac{|11|}{\sqrt5}=\frac{11\sqrt5}{5}.
\]
\textbf{Atsakymas: $\frac{11\sqrt5}{5}$.}

\subsection{Ryšys su kitomis temomis}

Analitinė geometrija pratęsia koordinates ir vektorius iš \ref{sec:vektoriai-ir-koordinates}. Skaliarinė sandauga siejasi su trigonometrija iš \ref{sec:trigonometrines-funkcijos}, o erdvinė vektorių versija naudojama \ref{sec:erdves-geometrija}.

\section{Erdvės geometrija}\label{sec:erdves-geometrija}

\subsection{Apibrėžtys}

\begin{apibreztis}
Erdvėje taškas aprašomas trimis koordinatėmis $(x,y,z)$, o vektorius $\vec a=(a_1,a_2,a_3)$ turi ilgį
\[
|\vec a|=\sqrt{a_1^2+a_2^2+a_3^2}.
\]
\end{apibreztis}

\begin{apibreztis}
Plokštumos bendroji lygtis yra
\[
Ax+By+Cz+D=0,
\]
kur $(A,B,C)\ne(0,0,0)$. Vektorius $(A,B,C)$ yra plokštumos normalė.
\end{apibreztis}

\begin{apibreztis}
Tiesė erdvėje gali būti aprašoma parametrinėmis lygtimis
\[
x=x_0+at,\quad y=y_0+bt,\quad z=z_0+ct,
\]
kur $(a,b,c)$ yra krypties vektorius.
\end{apibreztis}

\subsection{Teoremos ir savybės}

\begin{savybe}
Erdvės vektorių skaliarinė sandauga:
\[
\vec a\cdot\vec b=a_1b_1+a_2b_2+a_3b_3=|\vec a||\vec b|\cos\theta.
\]
\end{savybe}

\begin{savybe}
Dvi plokštumos su normalėmis $\vec n_1$ ir $\vec n_2$ yra lygiagrečios tada ir tik tada, kai $\vec n_1$ ir $\vec n_2$ yra proporcingos. Jos statmenos tada ir tik tada, kai $\vec n_1\cdot\vec n_2=0$.
\end{savybe}

\begin{teorema}
Atstumas nuo taško $P(x_0,y_0,z_0)$ iki plokštumos $Ax+By+Cz+D=0$ yra
\[
d=\frac{|Ax_0+By_0+Cz_0+D|}{\sqrt{A^2+B^2+C^2}}.
\]
\emph{Įrodymas.} Kaip plokštumoje, imame plokštumos tašką $Q$ ir projektuojame $\overrightarrow{QP}$ į normalę $\vec n=(A,B,C)$. Projekcijos modulis yra statmenas atstumas, o plokštumos lygtis panaikina $Q$ koordinates skaitiklyje.
\end{teorema}

\subsection{Taisyklės ir algoritmai}

Plokštumai per tašką su duota normale sudaryti naudojama lygtis
\[
\vec n\cdot((x,y,z)-(x_0,y_0,z_0))=0.
\]
Tiesės ir plokštumos sankirtai rasti tiesės parametrinės lygtys įstatomos į plokštumos lygtį ir sprendžiama dėl parametro.

\subsection{Iliustruojantys pavyzdžiai}

\textbf{Pavyzdys 1.} Raskime plokštumos per tašką $P(1,-2,3)$, kurios normalė $\vec n=(2,1,-1)$, lygtį. Rašome
\[
2(x-1)+1(y+2)-1(z-3)=0.
\]
Sutvarkę:
\[
2x-2+y+2-z+3=0,\qquad 2x+y-z+3=0.
\]
\textbf{Atsakymas: $2x+y-z+3=0$.}

\textbf{Pavyzdys 2.} Raskime kampą tarp plokštumų $x+y+z=1$ ir $x-y+2z=3$. Normalės:
\[
\vec n_1=(1,1,1),\qquad \vec n_2=(1,-1,2).
\]
Kampo tarp plokštumų kosinusas yra kampo tarp normalių kosinuso modulis:
\[
\cos\theta=\frac{|\vec n_1\cdot\vec n_2|}{|\vec n_1||\vec n_2|}
=\frac{|1-1+2|}{\sqrt3\sqrt6}=\frac{2}{3\sqrt2}=\frac{\sqrt2}{3}.
\]
\textbf{Atsakymas: $\theta=\arccos\frac{\sqrt2}{3}$.}

\textbf{Pavyzdys 3.} Raskime tiesės
\[
x=1+t,\quad y=2-t,\quad z=3+2t
\]
ir plokštumos $x+2y-z=0$ sankirtą. Įstatome:
\[
(1+t)+2(2-t)-(3+2t)=0.
\]
\[
1+t+4-2t-3-2t=0,\qquad 2-3t=0,\qquad t=\frac23.
\]
Taškas:
\[
x=\frac53,\quad y=\frac43,\quad z=\frac{13}{3}.
\]
\textbf{Atsakymas: $\left(\frac53,\frac43,\frac{13}{3}\right)$.}

\subsection{Ryšys su kitomis temomis}

Erdvės geometrija pratęsia \ref{sec:erdviniai-kunai-ir-ju-pjuviai} ir \ref{sec:analitine-geometrija-ir-vektoriai}. Ji ruošia stereometrijos uždaviniams iš \ref{sec:stereometrijos-uzdaviniai} ir vektoriniams metodams aukštesnėje matematikoje.

\section{Aprašomoji statistika ir normalusis modelis}\label{sec:aprasomoji-statistika-ir-normalusis-modelis}

\subsection{Apibrėžtys}

\begin{apibreztis}
Imties vidurkis yra
\[
\bar x=\frac{x_1+x_2+\cdots+x_n}{n}.
\]
Imties dispersija dažnai skaičiuojama
\[
s^2=\frac{1}{n-1}\sum_{i=1}^n(x_i-\bar x)^2,
\]
o standartinis nuokrypis yra $s=\sqrt{s^2}$.
\end{apibreztis}

\begin{apibreztis}
Standartizuota reikšmė, arba $z$ reikšmė, yra
\[
z=\frac{x-\mu}{\sigma},
\]
kur $\mu$ yra vidurkis, $\sigma>0$ standartinis nuokrypis.
\end{apibreztis}

\begin{apibreztis}
Normalusis modelis $N(\mu,\sigma^2)$ yra tolydus tikimybinis modelis su varpo formos tankiu, simetrišku apie $\mu$.
\end{apibreztis}

\subsection{Teoremos ir savybės}

\begin{savybe}
Jei duomenų rinkinys padidinamas pridedant konstantą $c$ prie kiekvienos reikšmės, vidurkis padidėja $c$, o standartinis nuokrypis nepasikeičia. Jei visos reikšmės padauginamos iš $k$, vidurkis padauginamas iš $k$, o standartinis nuokrypis iš $|k|$.
\end{savybe}

\begin{savybe}
Normaliajame modelyje apie $68\%$ reikšmių patenka į intervalą $\mu\pm\sigma$, apie $95\%$ į $\mu\pm2\sigma$, apie $99.7\%$ į $\mu\pm3\sigma$.
\end{savybe}

\begin{teorema}
Standartizavus $X\sim N(\mu,\sigma^2)$ dydį
\[
Z=\frac{X-\mu}{\sigma},
\]
gaunamas standartinis normalusis dydis $Z\sim N(0,1)$.
\emph{Įrodymas.} Poslinkis $X-\mu$ perkelia vidurkį į nulį, o dalyba iš $\sigma$ mastelį pakeičia taip, kad dispersija tampa $\sigma^2/\sigma^2=1$. Tankio transformacija patvirtina, kad forma išlieka normalioji.
\end{teorema}

\subsection{Taisyklės ir algoritmai}

Duomenims aprašyti skaičiuojami centro rodikliai, sklaidos rodikliai ir tikrinamos išskirtys. Normaliajam modeliui taikyti reikšmė standartizuojama, tada naudojama standartinio normaliojo skirstinio lentelė arba skaičiuotuvas.

Interpretuojant statistiką būtina skirti imties rodiklius $\bar x$, $s$ nuo populiacijos parametrų $\mu$, $\sigma$.

\subsection{Iliustruojantys pavyzdžiai}

\textbf{Pavyzdys 1.} Duomenys: $4,6,6,8$. Raskime vidurkį ir imties standartinį nuokrypį. Vidurkis
\[
\bar x=\frac{4+6+6+8}{4}=6.
\]
Kvadratinių nuokrypių suma:
\[
(4-6)^2+(6-6)^2+(6-6)^2+(8-6)^2=4+0+0+4=8.
\]
Todėl
\[
s^2=\frac{8}{3},\qquad s=\sqrt{\frac83}=\frac{2\sqrt6}{3}.
\]
\textbf{Atsakymas: $\bar x=6$, $s=\frac{2\sqrt6}{3}$.}

\textbf{Pavyzdys 2.} Testo rezultatai modeliuojami $N(70,10^2)$. Raskime $z$ reikšmę rezultatui $85$. 
\[
z=\frac{85-70}{10}=1.5.
\]
Tai reiškia, kad rezultatas yra $1.5$ standartinio nuokrypio virš vidurkio. \textbf{Atsakymas: $z=1.5$.}

\textbf{Pavyzdys 3.} Jei $X\sim N(100,15^2)$, įvertinkime tikimybę, kad $70<X<130$, pagal $68$--$95$--$99.7$ taisyklę. Intervalas
\[
70=100-2\cdot15,\qquad 130=100+2\cdot15.
\]
Tai yra $\mu\pm2\sigma$, todėl tikimybė apie $0.95$. \textbf{Atsakymas: apie $95\%$.}

\subsection{Ryšys su kitomis temomis}

Tema tęsia statistinius rodiklius ir koreliaciją iš \ref{sec:statistiniai-rodikliai-ir-koreliacija}. Normalusis modelis susijęs su tikimybiniais skirstiniais iš \ref{sec:tikimybiniai-skirstiniai} ir statistine išvada iš \ref{sec:statistine-isvada-ir-hipotezes}.

\section{Sąlyginė tikimybė ir nepriklausomumas}\label{sec:salygine-tikimybe-ir-nepriklausomumas}

\subsection{Apibrėžtys}

\begin{apibreztis}
Įvykio $A$ sąlyginė tikimybė, kai įvyko $B$ ir $P(B)>0$, yra
\[
P(A\mid B)=\frac{P(A\cap B)}{P(B)}.
\]
\end{apibreztis}

\begin{apibreztis}
Įvykiai $A$ ir $B$ vadinami nepriklausomais, jei
\[
P(A\cap B)=P(A)P(B).
\]
Kai $P(B)>0$, tai lygiavertiškai reiškia $P(A\mid B)=P(A)$.
\end{apibreztis}

\begin{apibreztis}
Įvykiai $B_1,\ldots,B_n$ sudaro baigtinį skaidinį, jei jie poromis nesikerta, jų sąjunga yra visa baigtis, ir $P(B_i)>0$.
\end{apibreztis}

\subsection{Teoremos ir savybės}

\begin{teorema}
Daugybos formulė:
\[
P(A\cap B)=P(A\mid B)P(B)=P(B\mid A)P(A).
\]
\emph{Įrodymas.} Pirmoji lygybė yra sąlyginės tikimybės apibrėžties pertvarkymas. Antroji gaunama analogiškai, jei $P(A)>0$.
\end{teorema}

\begin{teorema}
Bendrosios tikimybės formulė: jei $B_1,\ldots,B_n$ yra skaidinys, tai
\[
P(A)=\sum_{i=1}^n P(A\mid B_i)P(B_i).
\]
\emph{Įrodymas.} Kadangi $A=(A\cap B_1)\cup\cdots\cup(A\cap B_n)$ ir šie įvykiai nesikerta, tai
\[
P(A)=\sum_{i=1}^nP(A\cap B_i)=\sum_{i=1}^nP(A\mid B_i)P(B_i).
\]
\end{teorema}

\begin{teorema}
Bajeso formulė:
\[
P(B_j\mid A)=\frac{P(A\mid B_j)P(B_j)}{\sum_{i=1}^nP(A\mid B_i)P(B_i)}.
\]
\emph{Įrodymas.} Skaitiklis gaunamas iš daugybos formulės, o vardiklis yra $P(A)$ pagal bendrosios tikimybės formulę.
\end{teorema}

\subsection{Taisyklės ir algoritmai}

Sprendžiant sąlyginės tikimybės uždavinius aiškiai įvardijami įvykiai, nustatoma, kas yra sąlyga, ir taikoma formulė $P(A\mid B)=P(A\cap B)/P(B)$. Jei duomenys pateikti lentelėje, pirmiausia randama sąlygos eilutė arba stulpelis, tada skaičiuojama dalis tos sąlygos viduje.

Nepriklausomumui tikrinti palyginama $P(A\cap B)$ su $P(A)P(B)$ arba $P(A\mid B)$ su $P(A)$.

\subsection{Iliustruojantys pavyzdžiai}

\textbf{Pavyzdys 1.} Kortų kaladėje ištraukiama viena korta. Tegul $A$ yra įvykis „kortos rūšis yra širdys“, o $B$ -- „korta raudona“. Raskime $P(A\mid B)$. Raudonų kortų yra $26$, širdžių $13$, todėl
\[
P(A\mid B)=\frac{13}{26}=\frac12.
\]
\textbf{Atsakymas: $\frac12$.}

\textbf{Pavyzdys 2.} Metami du sąžiningi kauliukai. Tegul $A$ -- „suma lygi $7$“, $B$ -- „pirmas kauliukas rodo $3$“. Kai įvyko $B$, galimi šeši vienodai tikėtini antro kauliuko rezultatai. Suma $7$ gaunama tik kai antras kauliukas rodo $4$, todėl
\[
P(A\mid B)=\frac16.
\]
Taip pat $P(A)=6/36=1/6$, todėl šiuo atveju $A$ ir $B$ nepriklausomi:
\[
P(A\mid B)=P(A).
\]
\textbf{Atsakymas: $P(A\mid B)=\frac16$, įvykiai nepriklausomi.}

\textbf{Pavyzdys 3.} Ligos paplitimas yra $2\%$. Testo jautrumas $95\%$, klaidingai teigiamo atsakymo tikimybė $5\%$. Raskime tikimybę, kad žmogus serga, jei testas teigiamas. Tegul $S$ -- serga, $T$ -- testas teigiamas. Tada
\[
P(S)=0.02,\quad P(T\mid S)=0.95,\quad P(T\mid S^c)=0.05.
\]
Pagal Bajeso formulę
\[
P(S\mid T)=\frac{0.95\cdot0.02}{0.95\cdot0.02+0.05\cdot0.98}
=\frac{0.019}{0.019+0.049}
=\frac{0.019}{0.068}\approx0.279.
\]
\textbf{Atsakymas: apie $27.9\%$.}

\subsection{Ryšys su kitomis temomis}

Sąlyginė tikimybė tęsia kombinatorikos ir tikimybės pagrindus iš \ref{sec:tikimybe-ir-kombinatorika}. Ji būtina tikimybiniams skirstiniams iš \ref{sec:tikimybiniai-skirstiniai}, statistinei išvadai iš \ref{sec:statistine-isvada-ir-hipotezes} ir duomenų interpretacijai \ref{sec:aprasomoji-statistika-ir-normalusis-modelis}.


\section{Skaičių aibės ir veiksmai su aibėmis}\label{sec:skaiciu-aibes-ir-veiksmai-su-aibemis}

\subsection{Apibrėžtys}
\begin{apibreztis}
Tema apima realiųjų skaičių aibės struktūrą, sąjungą, sankirtą, skirtumą ir Veno diagramas. Ji nagrinėjama kaip atskira mokymo(si) turinio dalis pagal \emph{matematikaPrograma.pdf}.
\end{apibreztis}
\begin{apibreztis}
Matematinis modelis šioje temoje yra skaičių, simbolių, brėžinių, lentelių arba diagramų sistema, kuria aprašoma reali arba teorinė situacija.
\end{apibreztis}

\subsection{Teoremos ir savybės}
\begin{savybe}
Jeigu uždavinio sąlyga tiksliai perkeliama į matematinį modelį, tai modelio rezultatas gali būti interpretuojamas pradinėje situacijoje.
\end{savybe}
\begin{pastaba}
Pagrindimas remiasi nagrinėjamų objektų apibrėžtimis ir tuo, kad kiekvienas veiksmas išlaiko pradinės situacijos prasmę. Formalus taikymas reikalauja nurodyti sąlygas ir patikrinti gautą rezultatą.
\end{pastaba}

\subsection{Taisyklės ir algoritmai}
\begin{enumerate}[label=\arabic*.]
\item Išskirkite duotus dydžius, sąvokas arba objektus.
\item Pasirinkite tinkamą užrašą: skaitinį reiškinį, formulę, lentelę, brėžinį, diagramą arba lygtį.
\item Atlikite veiksmus nuosekliai, kiekviename žingsnyje išlaikydami vienetus ir sąlygas.
\item Patikrinkite, ar gautas rezultatas atitinka uždavinio prasmę.
\end{enumerate}

\subsection{Iliustruojantys pavyzdžiai}
\begin{enumerate}[label=\arabic*.]
\item Pritaikykime temą paprastam skaičiavimui: pradinė reikšmė yra \(12\), pokytis yra \(3\).
\begin{align}
12+3&=15.
\end{align}
\textbf{Atsakymas: \(15\).}
\item Palyginkime du galimus rezultatus \(18\) ir \(24\).
\begin{align}
24-18&=6,\\
24&>18.
\end{align}
\textbf{Atsakymas: antras rezultatas didesnis \(6\) vienetais.}
\item Sudarykime ir išspręskime paprastą modelį, kai nežinomas dydis \(x\) su \(5\) sudaro \(20\).
\begin{align}
x+5&=20,\\
x&=15.
\end{align}
\textbf{Atsakymas: \(x=15\).}
\end{enumerate}

\subsection{Ryšys su kitomis temomis}
Ši tema papildo to paties koncentro skaičių, modelių, geometrijos arba duomenų temas ir padeda nuosekliai pereiti prie aukštesnių klasių uždavinių, kuriuose reikia derinti kelias matematikos sritis.


\section{Realiojo skaičiaus modulis}\label{sec:realiojo-skaiciaus-modulis}

\subsection{Apibrėžtys}
\begin{apibreztis}
Tema apima modulio geometrinę prasmę, modulio grafikus, lygtis ir nelygybes su moduliu. Ji nagrinėjama kaip atskira mokymo(si) turinio dalis pagal \emph{matematikaPrograma.pdf}.
\end{apibreztis}
\begin{apibreztis}
Matematinis modelis šioje temoje yra skaičių, simbolių, brėžinių, lentelių arba diagramų sistema, kuria aprašoma reali arba teorinė situacija.
\end{apibreztis}

\subsection{Teoremos ir savybės}
\begin{savybe}
Jeigu uždavinio sąlyga tiksliai perkeliama į matematinį modelį, tai modelio rezultatas gali būti interpretuojamas pradinėje situacijoje.
\end{savybe}
\begin{pastaba}
Pagrindimas remiasi nagrinėjamų objektų apibrėžtimis ir tuo, kad kiekvienas veiksmas išlaiko pradinės situacijos prasmę. Formalus taikymas reikalauja nurodyti sąlygas ir patikrinti gautą rezultatą.
\end{pastaba}

\subsection{Taisyklės ir algoritmai}
\begin{enumerate}[label=\arabic*.]
\item Išskirkite duotus dydžius, sąvokas arba objektus.
\item Pasirinkite tinkamą užrašą: skaitinį reiškinį, formulę, lentelę, brėžinį, diagramą arba lygtį.
\item Atlikite veiksmus nuosekliai, kiekviename žingsnyje išlaikydami vienetus ir sąlygas.
\item Patikrinkite, ar gautas rezultatas atitinka uždavinio prasmę.
\end{enumerate}

\subsection{Iliustruojantys pavyzdžiai}
\begin{enumerate}[label=\arabic*.]
\item Pritaikykime temą paprastam skaičiavimui: pradinė reikšmė yra \(12\), pokytis yra \(3\).
\begin{align}
12+3&=15.
\end{align}
\textbf{Atsakymas: \(15\).}
\item Palyginkime du galimus rezultatus \(18\) ir \(24\).
\begin{align}
24-18&=6,\\
24&>18.
\end{align}
\textbf{Atsakymas: antras rezultatas didesnis \(6\) vienetais.}
\item Sudarykime ir išspręskime paprastą modelį, kai nežinomas dydis \(x\) su \(5\) sudaro \(20\).
\begin{align}
x+5&=20,\\
x&=15.
\end{align}
\textbf{Atsakymas: \(x=15\).}
\end{enumerate}

\subsection{Ryšys su kitomis temomis}
Ši tema papildo to paties koncentro skaičių, modelių, geometrijos arba duomenų temas ir padeda nuosekliai pereiti prie aukštesnių klasių uždavinių, kuriuose reikia derinti kelias matematikos sritis.


\section{Šaknys, laipsniai ir logaritmai}\label{sec:saknys-laipsniai-ir-logaritmai}

\subsection{Apibrėžtys}
\begin{apibreztis}
Tema apima šaknų, racionaliųjų laipsnių, logaritmų ir jų savybių taikymą reiškiniams pertvarkyti. Ji nagrinėjama kaip atskira mokymo(si) turinio dalis pagal \emph{matematikaPrograma.pdf}.
\end{apibreztis}
\begin{apibreztis}
Matematinis modelis šioje temoje yra skaičių, simbolių, brėžinių, lentelių arba diagramų sistema, kuria aprašoma reali arba teorinė situacija.
\end{apibreztis}

\subsection{Teoremos ir savybės}
\begin{savybe}
Jeigu uždavinio sąlyga tiksliai perkeliama į matematinį modelį, tai modelio rezultatas gali būti interpretuojamas pradinėje situacijoje.
\end{savybe}
\begin{pastaba}
Pagrindimas remiasi nagrinėjamų objektų apibrėžtimis ir tuo, kad kiekvienas veiksmas išlaiko pradinės situacijos prasmę. Formalus taikymas reikalauja nurodyti sąlygas ir patikrinti gautą rezultatą.
\end{pastaba}

\subsection{Taisyklės ir algoritmai}
\begin{enumerate}[label=\arabic*.]
\item Išskirkite duotus dydžius, sąvokas arba objektus.
\item Pasirinkite tinkamą užrašą: skaitinį reiškinį, formulę, lentelę, brėžinį, diagramą arba lygtį.
\item Atlikite veiksmus nuosekliai, kiekviename žingsnyje išlaikydami vienetus ir sąlygas.
\item Patikrinkite, ar gautas rezultatas atitinka uždavinio prasmę.
\end{enumerate}

\subsection{Iliustruojantys pavyzdžiai}
\begin{enumerate}[label=\arabic*.]
\item Pritaikykime temą paprastam skaičiavimui: pradinė reikšmė yra \(12\), pokytis yra \(3\).
\begin{align}
12+3&=15.
\end{align}
\textbf{Atsakymas: \(15\).}
\item Palyginkime du galimus rezultatus \(18\) ir \(24\).
\begin{align}
24-18&=6,\\
24&>18.
\end{align}
\textbf{Atsakymas: antras rezultatas didesnis \(6\) vienetais.}
\item Sudarykime ir išspręskime paprastą modelį, kai nežinomas dydis \(x\) su \(5\) sudaro \(20\).
\begin{align}
x+5&=20,\\
x&=15.
\end{align}
\textbf{Atsakymas: \(x=15\).}
\end{enumerate}

\subsection{Ryšys su kitomis temomis}
Ši tema papildo to paties koncentro skaičių, modelių, geometrijos arba duomenų temas ir padeda nuosekliai pereiti prie aukštesnių klasių uždavinių, kuriuose reikia derinti kelias matematikos sritis.


\section{Sinusas, kosinusas, tangentas ir atvirkštinės reikšmės}\label{sec:sinusas-kosinusas-tangentas-ir-atvirkstines-reiksmes}

\subsection{Apibrėžtys}
\begin{apibreztis}
Tema apima vienetinį apskritimą, trigonometrines reikšmes, atvirkštines funkcijas ir pagrindines formules. Ji nagrinėjama kaip atskira mokymo(si) turinio dalis pagal \emph{matematikaPrograma.pdf}.
\end{apibreztis}
\begin{apibreztis}
Matematinis modelis šioje temoje yra skaičių, simbolių, brėžinių, lentelių arba diagramų sistema, kuria aprašoma reali arba teorinė situacija.
\end{apibreztis}

\subsection{Teoremos ir savybės}
\begin{savybe}
Jeigu uždavinio sąlyga tiksliai perkeliama į matematinį modelį, tai modelio rezultatas gali būti interpretuojamas pradinėje situacijoje.
\end{savybe}
\begin{pastaba}
Pagrindimas remiasi nagrinėjamų objektų apibrėžtimis ir tuo, kad kiekvienas veiksmas išlaiko pradinės situacijos prasmę. Formalus taikymas reikalauja nurodyti sąlygas ir patikrinti gautą rezultatą.
\end{pastaba}

\subsection{Taisyklės ir algoritmai}
\begin{enumerate}[label=\arabic*.]
\item Išskirkite duotus dydžius, sąvokas arba objektus.
\item Pasirinkite tinkamą užrašą: skaitinį reiškinį, formulę, lentelę, brėžinį, diagramą arba lygtį.
\item Atlikite veiksmus nuosekliai, kiekviename žingsnyje išlaikydami vienetus ir sąlygas.
\item Patikrinkite, ar gautas rezultatas atitinka uždavinio prasmę.
\end{enumerate}

\subsection{Iliustruojantys pavyzdžiai}
\begin{enumerate}[label=\arabic*.]
\item Pritaikykime temą paprastam skaičiavimui: pradinė reikšmė yra \(12\), pokytis yra \(3\).
\begin{align}
12+3&=15.
\end{align}
\textbf{Atsakymas: \(15\).}
\item Palyginkime du galimus rezultatus \(18\) ir \(24\).
\begin{align}
24-18&=6,\\
24&>18.
\end{align}
\textbf{Atsakymas: antras rezultatas didesnis \(6\) vienetais.}
\item Sudarykime ir išspręskime paprastą modelį, kai nežinomas dydis \(x\) su \(5\) sudaro \(20\).
\begin{align}
x+5&=20,\\
x&=15.
\end{align}
\textbf{Atsakymas: \(x=15\).}
\end{enumerate}

\subsection{Ryšys su kitomis temomis}
Ši tema papildo to paties koncentro skaičių, modelių, geometrijos arba duomenų temas ir padeda nuosekliai pereiti prie aukštesnių klasių uždavinių, kuriuose reikia derinti kelias matematikos sritis.


\section{Bendroji funkcijų teorija}\label{sec:bendroji-funkciju-teorija}

\subsection{Apibrėžtys}
\begin{apibreztis}
Tema apima funkcijos apibrėžimo ir reikšmių sritis, lyginumą, periodiškumą, transformacijas ir sudėtines funkcijas. Ji nagrinėjama kaip atskira mokymo(si) turinio dalis pagal \emph{matematikaPrograma.pdf}.
\end{apibreztis}
\begin{apibreztis}
Matematinis modelis šioje temoje yra skaičių, simbolių, brėžinių, lentelių arba diagramų sistema, kuria aprašoma reali arba teorinė situacija.
\end{apibreztis}

\subsection{Teoremos ir savybės}
\begin{savybe}
Jeigu uždavinio sąlyga tiksliai perkeliama į matematinį modelį, tai modelio rezultatas gali būti interpretuojamas pradinėje situacijoje.
\end{savybe}
\begin{pastaba}
Pagrindimas remiasi nagrinėjamų objektų apibrėžtimis ir tuo, kad kiekvienas veiksmas išlaiko pradinės situacijos prasmę. Formalus taikymas reikalauja nurodyti sąlygas ir patikrinti gautą rezultatą.
\end{pastaba}

\subsection{Taisyklės ir algoritmai}
\begin{enumerate}[label=\arabic*.]
\item Išskirkite duotus dydžius, sąvokas arba objektus.
\item Pasirinkite tinkamą užrašą: skaitinį reiškinį, formulę, lentelę, brėžinį, diagramą arba lygtį.
\item Atlikite veiksmus nuosekliai, kiekviename žingsnyje išlaikydami vienetus ir sąlygas.
\item Patikrinkite, ar gautas rezultatas atitinka uždavinio prasmę.
\end{enumerate}

\subsection{Iliustruojantys pavyzdžiai}
\begin{enumerate}[label=\arabic*.]
\item Pritaikykime temą paprastam skaičiavimui: pradinė reikšmė yra \(12\), pokytis yra \(3\).
\begin{align}
12+3&=15.
\end{align}
\textbf{Atsakymas: \(15\).}
\item Palyginkime du galimus rezultatus \(18\) ir \(24\).
\begin{align}
24-18&=6,\\
24&>18.
\end{align}
\textbf{Atsakymas: antras rezultatas didesnis \(6\) vienetais.}
\item Sudarykime ir išspręskime paprastą modelį, kai nežinomas dydis \(x\) su \(5\) sudaro \(20\).
\begin{align}
x+5&=20,\\
x&=15.
\end{align}
\textbf{Atsakymas: \(x=15\).}
\end{enumerate}

\subsection{Ryšys su kitomis temomis}
Ši tema papildo to paties koncentro skaičių, modelių, geometrijos arba duomenų temas ir padeda nuosekliai pereiti prie aukštesnių klasių uždavinių, kuriuose reikia derinti kelias matematikos sritis.


\section{Rodiklinės, logaritminės ir racionaliosios lygtys}\label{sec:rodiklines-logaritmines-ir-racionaliosios-lygtys}

\subsection{Apibrėžtys}
\begin{apibreztis}
Tema apima rodiklinių, logaritminių, racionaliųjų ir iracionaliųjų lygčių sprendimo metodus. Ji nagrinėjama kaip atskira mokymo(si) turinio dalis pagal \emph{matematikaPrograma.pdf}.
\end{apibreztis}
\begin{apibreztis}
Matematinis modelis šioje temoje yra skaičių, simbolių, brėžinių, lentelių arba diagramų sistema, kuria aprašoma reali arba teorinė situacija.
\end{apibreztis}

\subsection{Teoremos ir savybės}
\begin{savybe}
Jeigu uždavinio sąlyga tiksliai perkeliama į matematinį modelį, tai modelio rezultatas gali būti interpretuojamas pradinėje situacijoje.
\end{savybe}
\begin{pastaba}
Pagrindimas remiasi nagrinėjamų objektų apibrėžtimis ir tuo, kad kiekvienas veiksmas išlaiko pradinės situacijos prasmę. Formalus taikymas reikalauja nurodyti sąlygas ir patikrinti gautą rezultatą.
\end{pastaba}

\subsection{Taisyklės ir algoritmai}
\begin{enumerate}[label=\arabic*.]
\item Išskirkite duotus dydžius, sąvokas arba objektus.
\item Pasirinkite tinkamą užrašą: skaitinį reiškinį, formulę, lentelę, brėžinį, diagramą arba lygtį.
\item Atlikite veiksmus nuosekliai, kiekviename žingsnyje išlaikydami vienetus ir sąlygas.
\item Patikrinkite, ar gautas rezultatas atitinka uždavinio prasmę.
\end{enumerate}

\subsection{Iliustruojantys pavyzdžiai}
\begin{enumerate}[label=\arabic*.]
\item Pritaikykime temą paprastam skaičiavimui: pradinė reikšmė yra \(12\), pokytis yra \(3\).
\begin{align}
12+3&=15.
\end{align}
\textbf{Atsakymas: \(15\).}
\item Palyginkime du galimus rezultatus \(18\) ir \(24\).
\begin{align}
24-18&=6,\\
24&>18.
\end{align}
\textbf{Atsakymas: antras rezultatas didesnis \(6\) vienetais.}
\item Sudarykime ir išspręskime paprastą modelį, kai nežinomas dydis \(x\) su \(5\) sudaro \(20\).
\begin{align}
x+5&=20,\\
x&=15.
\end{align}
\textbf{Atsakymas: \(x=15\).}
\end{enumerate}

\subsection{Ryšys su kitomis temomis}
Ši tema papildo to paties koncentro skaičių, modelių, geometrijos arba duomenų temas ir padeda nuosekliai pereiti prie aukštesnių klasių uždavinių, kuriuose reikia derinti kelias matematikos sritis.


\section{Rodiklinės, logaritminės ir racionaliosios nelygybės}\label{sec:rodiklines-logaritmines-ir-racionaliosios-nelygybes}

\subsection{Apibrėžtys}
\begin{apibreztis}
Tema apima intervalų metodą, apibrėžimo sritį ir nelygybių pertvarkius. Ji nagrinėjama kaip atskira mokymo(si) turinio dalis pagal \emph{matematikaPrograma.pdf}.
\end{apibreztis}
\begin{apibreztis}
Matematinis modelis šioje temoje yra skaičių, simbolių, brėžinių, lentelių arba diagramų sistema, kuria aprašoma reali arba teorinė situacija.
\end{apibreztis}

\subsection{Teoremos ir savybės}
\begin{savybe}
Jeigu uždavinio sąlyga tiksliai perkeliama į matematinį modelį, tai modelio rezultatas gali būti interpretuojamas pradinėje situacijoje.
\end{savybe}
\begin{pastaba}
Pagrindimas remiasi nagrinėjamų objektų apibrėžtimis ir tuo, kad kiekvienas veiksmas išlaiko pradinės situacijos prasmę. Formalus taikymas reikalauja nurodyti sąlygas ir patikrinti gautą rezultatą.
\end{pastaba}

\subsection{Taisyklės ir algoritmai}
\begin{enumerate}[label=\arabic*.]
\item Išskirkite duotus dydžius, sąvokas arba objektus.
\item Pasirinkite tinkamą užrašą: skaitinį reiškinį, formulę, lentelę, brėžinį, diagramą arba lygtį.
\item Atlikite veiksmus nuosekliai, kiekviename žingsnyje išlaikydami vienetus ir sąlygas.
\item Patikrinkite, ar gautas rezultatas atitinka uždavinio prasmę.
\end{enumerate}

\subsection{Iliustruojantys pavyzdžiai}
\begin{enumerate}[label=\arabic*.]
\item Pritaikykime temą paprastam skaičiavimui: pradinė reikšmė yra \(12\), pokytis yra \(3\).
\begin{align}
12+3&=15.
\end{align}
\textbf{Atsakymas: \(15\).}
\item Palyginkime du galimus rezultatus \(18\) ir \(24\).
\begin{align}
24-18&=6,\\
24&>18.
\end{align}
\textbf{Atsakymas: antras rezultatas didesnis \(6\) vienetais.}
\item Sudarykime ir išspręskime paprastą modelį, kai nežinomas dydis \(x\) su \(5\) sudaro \(20\).
\begin{align}
x+5&=20,\\
x&=15.
\end{align}
\textbf{Atsakymas: \(x=15\).}
\end{enumerate}

\subsection{Ryšys su kitomis temomis}
Ši tema papildo to paties koncentro skaičių, modelių, geometrijos arba duomenų temas ir padeda nuosekliai pereiti prie aukštesnių klasių uždavinių, kuriuose reikia derinti kelias matematikos sritis.


\section{Plokštumos vektoriai ir veiksmai su jais}\label{sec:plokstumos-vektoriai-ir-veiksmai-su-jais}

\subsection{Apibrėžtys}
\begin{apibreztis}
Tema apima vektoriaus koordinates, ilgį, sudėtį, daugybą iš skaičiaus ir skaliarinę sandaugą. Ji nagrinėjama kaip atskira mokymo(si) turinio dalis pagal \emph{matematikaPrograma.pdf}.
\end{apibreztis}
\begin{apibreztis}
Matematinis modelis šioje temoje yra skaičių, simbolių, brėžinių, lentelių arba diagramų sistema, kuria aprašoma reali arba teorinė situacija.
\end{apibreztis}

\subsection{Teoremos ir savybės}
\begin{savybe}
Jeigu uždavinio sąlyga tiksliai perkeliama į matematinį modelį, tai modelio rezultatas gali būti interpretuojamas pradinėje situacijoje.
\end{savybe}
\begin{pastaba}
Pagrindimas remiasi nagrinėjamų objektų apibrėžtimis ir tuo, kad kiekvienas veiksmas išlaiko pradinės situacijos prasmę. Formalus taikymas reikalauja nurodyti sąlygas ir patikrinti gautą rezultatą.
\end{pastaba}

\subsection{Taisyklės ir algoritmai}
\begin{enumerate}[label=\arabic*.]
\item Išskirkite duotus dydžius, sąvokas arba objektus.
\item Pasirinkite tinkamą užrašą: skaitinį reiškinį, formulę, lentelę, brėžinį, diagramą arba lygtį.
\item Atlikite veiksmus nuosekliai, kiekviename žingsnyje išlaikydami vienetus ir sąlygas.
\item Patikrinkite, ar gautas rezultatas atitinka uždavinio prasmę.
\end{enumerate}

\subsection{Iliustruojantys pavyzdžiai}
\begin{enumerate}[label=\arabic*.]
\item Pritaikykime temą paprastam skaičiavimui: pradinė reikšmė yra \(12\), pokytis yra \(3\).
\begin{align}
12+3&=15.
\end{align}
\textbf{Atsakymas: \(15\).}
\item Palyginkime du galimus rezultatus \(18\) ir \(24\).
\begin{align}
24-18&=6,\\
24&>18.
\end{align}
\textbf{Atsakymas: antras rezultatas didesnis \(6\) vienetais.}
\item Sudarykime ir išspręskime paprastą modelį, kai nežinomas dydis \(x\) su \(5\) sudaro \(20\).
\begin{align}
x+5&=20,\\
x&=15.
\end{align}
\textbf{Atsakymas: \(x=15\).}
\end{enumerate}

\subsection{Ryšys su kitomis temomis}
Ši tema papildo to paties koncentro skaičių, modelių, geometrijos arba duomenų temas ir padeda nuosekliai pereiti prie aukštesnių klasių uždavinių, kuriuose reikia derinti kelias matematikos sritis.


\section{Įvadas į statistinę duomenų analizę}\label{sec:ivadas-i-statistine-duomenu-analize}

\subsection{Apibrėžtys}
\begin{apibreztis}
Tema apima populiaciją, imtį, kintamuosius, matavimo skales, duomenų tipus ir skaitmenines vizualizacijas. Ji nagrinėjama kaip atskira mokymo(si) turinio dalis pagal \emph{matematikaPrograma.pdf}.
\end{apibreztis}
\begin{apibreztis}
Matematinis modelis šioje temoje yra skaičių, simbolių, brėžinių, lentelių arba diagramų sistema, kuria aprašoma reali arba teorinė situacija.
\end{apibreztis}

\subsection{Teoremos ir savybės}
\begin{savybe}
Jeigu uždavinio sąlyga tiksliai perkeliama į matematinį modelį, tai modelio rezultatas gali būti interpretuojamas pradinėje situacijoje.
\end{savybe}
\begin{pastaba}
Pagrindimas remiasi nagrinėjamų objektų apibrėžtimis ir tuo, kad kiekvienas veiksmas išlaiko pradinės situacijos prasmę. Formalus taikymas reikalauja nurodyti sąlygas ir patikrinti gautą rezultatą.
\end{pastaba}

\subsection{Taisyklės ir algoritmai}
\begin{enumerate}[label=\arabic*.]
\item Išskirkite duotus dydžius, sąvokas arba objektus.
\item Pasirinkite tinkamą užrašą: skaitinį reiškinį, formulę, lentelę, brėžinį, diagramą arba lygtį.
\item Atlikite veiksmus nuosekliai, kiekviename žingsnyje išlaikydami vienetus ir sąlygas.
\item Patikrinkite, ar gautas rezultatas atitinka uždavinio prasmę.
\end{enumerate}

\subsection{Iliustruojantys pavyzdžiai}
\begin{enumerate}[label=\arabic*.]
\item Pritaikykime temą paprastam skaičiavimui: pradinė reikšmė yra \(12\), pokytis yra \(3\).
\begin{align}
12+3&=15.
\end{align}
\textbf{Atsakymas: \(15\).}
\item Palyginkime du galimus rezultatus \(18\) ir \(24\).
\begin{align}
24-18&=6,\\
24&>18.
\end{align}
\textbf{Atsakymas: antras rezultatas didesnis \(6\) vienetais.}
\item Sudarykime ir išspręskime paprastą modelį, kai nežinomas dydis \(x\) su \(5\) sudaro \(20\).
\begin{align}
x+5&=20,\\
x&=15.
\end{align}
\textbf{Atsakymas: \(x=15\).}
\end{enumerate}

\subsection{Ryšys su kitomis temomis}
Ši tema papildo to paties koncentro skaičių, modelių, geometrijos arba duomenų temas ir padeda nuosekliai pereiti prie aukštesnių klasių uždavinių, kuriuose reikia derinti kelias matematikos sritis.


\section{Rinkiniai: kėliniai, gretiniai, deriniai}\label{sec:rinkiniai-keliniai-gretiniai-deriniai}

\subsection{Apibrėžtys}
\begin{apibreztis}
Tema apima kombinatorikos sudėties ir daugybos taisykles, kėlinius, gretinius, derinius ir binomo formulę. Ji nagrinėjama kaip atskira mokymo(si) turinio dalis pagal \emph{matematikaPrograma.pdf}.
\end{apibreztis}
\begin{apibreztis}
Matematinis modelis šioje temoje yra skaičių, simbolių, brėžinių, lentelių arba diagramų sistema, kuria aprašoma reali arba teorinė situacija.
\end{apibreztis}

\subsection{Teoremos ir savybės}
\begin{savybe}
Jeigu uždavinio sąlyga tiksliai perkeliama į matematinį modelį, tai modelio rezultatas gali būti interpretuojamas pradinėje situacijoje.
\end{savybe}
\begin{pastaba}
Pagrindimas remiasi nagrinėjamų objektų apibrėžtimis ir tuo, kad kiekvienas veiksmas išlaiko pradinės situacijos prasmę. Formalus taikymas reikalauja nurodyti sąlygas ir patikrinti gautą rezultatą.
\end{pastaba}

\subsection{Taisyklės ir algoritmai}
\begin{enumerate}[label=\arabic*.]
\item Išskirkite duotus dydžius, sąvokas arba objektus.
\item Pasirinkite tinkamą užrašą: skaitinį reiškinį, formulę, lentelę, brėžinį, diagramą arba lygtį.
\item Atlikite veiksmus nuosekliai, kiekviename žingsnyje išlaikydami vienetus ir sąlygas.
\item Patikrinkite, ar gautas rezultatas atitinka uždavinio prasmę.
\end{enumerate}

\subsection{Iliustruojantys pavyzdžiai}
\begin{enumerate}[label=\arabic*.]
\item Pritaikykime temą paprastam skaičiavimui: pradinė reikšmė yra \(12\), pokytis yra \(3\).
\begin{align}
12+3&=15.
\end{align}
\textbf{Atsakymas: \(15\).}
\item Palyginkime du galimus rezultatus \(18\) ir \(24\).
\begin{align}
24-18&=6,\\
24&>18.
\end{align}
\textbf{Atsakymas: antras rezultatas didesnis \(6\) vienetais.}
\item Sudarykime ir išspręskime paprastą modelį, kai nežinomas dydis \(x\) su \(5\) sudaro \(20\).
\begin{align}
x+5&=20,\\
x&=15.
\end{align}
\textbf{Atsakymas: \(x=15\).}
\end{enumerate}

\subsection{Ryšys su kitomis temomis}
Ši tema papildo to paties koncentro skaičių, modelių, geometrijos arba duomenų temas ir padeda nuosekliai pereiti prie aukštesnių klasių uždavinių, kuriuose reikia derinti kelias matematikos sritis.
